\documentclass{article}
\usepackage[T1]{fontenc}
\usepackage[utf8]{inputenc}
\usepackage{url}
\usepackage{listings}
\usepackage{tikz}
\usepackage{float}

\lstset{
  basicstyle=\ttfamily\footnotesize,
  breaklines=true,
  breakatwhitespace=false,
  columns=fullflexible,
  keepspaces=true
}

\tikzset{
  vertex/.style={circle, draw, minimum size=8mm, inner sep=0pt},
  frontier/.style={circle, draw, fill=blue!25, minimum size=8mm, inner sep=0pt},
  visited/.style={circle, draw, fill=gray!25, minimum size=8mm, inner sep=0pt},
  undiscovered/.style={circle, draw, fill=white, minimum size=8mm, inner sep=0pt},
  treeedge/.style={very thick, blue!70!black},
  graphedge/.style={gray!70}
}

\begin{document}
\section{Refatoracao do kernel Graph500}
Um dos objetivos deste trabalho era reimplementar a versao de referencia do
benchmark Graph500 utilizando NVSHMEM e mantendo o algoritimo original. 
A implementacao segue um modelo hibrido, o host controla a sincronizacao 
entre fronteiras da BFS, enquanto um kernel 
CUDA executa a expansao da fronteira atual e utiliza operacoes NVSHMEM para atualizar 
dados remotos diretamente entre PEs. Assim, a CPU ainda participa da coordenacao global 
da busca, mas deixa de ser necessaria para cada descoberta remota de vertice.

\section{Estrutura do Graph500}
Na versao de referencia do Graph500, o fluxo principal do benchmark e separado
da implementacao da travessia por uma pequena interface. O programa principal
gera o grafo, escolhe as raizes da BFS, chama uma rotina para construir a
estrutura interna do grafo e, para cada raiz, chama a funcao de BFS.

A primeira etapa e a conversao do grafo gerado em tuplas para uma representacao
CSR distribuida. O Graph500 distribui os vertices de forma ciclica entre os
ranks. O dono de um vertice global \texttt{v} e calculado por \texttt{v mod P},
onde \texttt{P} e o numero de PEs, e o indice local e calculado por \texttt{v / P}.

Durante o algoritmo da BFS enquanto a fronteira global nao estiver vazia, 
o host lanca o kernel de expansao
da fronteira. Dentro desse kernel, cada thread processa vertices da fronteira
local. Para cada vertice, a thread percorre sua lista de adjacencia no CSR,
calcula o PE dono do vizinho e tenta marcar esse vizinho como descoberto.

Quando uma descoberta e bem-sucedida, o kernel tambem insere o vertice na
proxima fronteira do PE que possui esse vertice. Para isso, usa uma operacao
atomica remota de incremento sobre o contador da proxima fronteira e, em
seguida, escreve o indice local do vertice descoberto na posicao obtida. Esse
padrao e o ponto central do uso de NVSHMEM na BFS, a thread que encontra uma
aresta remota nao precisa retornar ao host nem empacotar uma mensagem MPI. Ela
atualiza diretamente o estado remoto necessario para que o PE dono processe esse
vertice no proximo nivel.

Ao final da expansao, o kernel garante a conclusao das operacoes remotas, e o
host volta a sincronizar os PEs. O contador da proxima fronteira e copiado da
GPU para a CPU, os ranks executam uma reducao global para somar os tamanhos das
fronteiras locais, e a BFS continua apenas se essa soma for maior que zero.

Essa divisao entre host e device torna a implementacao mais simples de validar,
pois preserva a estrutura classica da BFS por fronteiras. Ao mesmo tempo, ela
mostra uma limitacao importante, embora as descobertas remotas ocorram no
device, a cada nivel ainda existem sincronizacoes com a CPU, copias de
contadores e reducoes globais.

\section{Geracao de grafo na GPU}
A geracao do grafo na implementacao refatorada foi movida para a GPU. O objetivo
e produzir a lista de arestas do Graph500 diretamente em memoria de dispositivo,
evitando que a CPU gere todas as tuplas e depois copie uma entrada grande para a
GPU. O gerador recebe um intervalo de arestas, esse
intervalo permite que cada PE gere apenas sua parte da lista global de arestas.
Assim, a geracao fica naturalmente paralela entre PEs e tambem paralela dentro
de cada GPU.

A implementacao tambem gera um peso associado a cada aresta. Esse peso e
produzido a partir do mesmo estado pseudoaleatorio usado na geracao da aresta,
permitindo experimentos com SSSP, que e outro kernel previsto pelo Graph500.

\section{BFS somente kernel}
No trabalho tambem foi feita uma implementacao que utiliza somente o kernel, sem
sincronizacao feita pelo host entre os niveis da BFS. O codigo segue a mesma
estrutura conceitual da versao com host: existem duas fronteiras, um vetor de
predecessores, contadores de fronteira e operacoes NVSHMEM para atualizar
vertices remotos. A diferenca principal e que o laco da BFS deixa de estar no
host e passa a estar dentro de um unico kernel CUDA.

A principal vantagem dessa abordagem e reduzir o custo fixo por nivel da BFS.
Em grafos de pequeno diametro, esse custo pode ser menos dominante, mas em
execucoes com muitos niveis a sincronizacao host-device repetida pode afetar o
tempo total. Ao manter o laco dentro do kernel, a implementacao evita multiplos
lancamentos de kernel, copias de contadores para a CPU e reducoes MPI entre
niveis. Isso aproxima a BFS do modelo de comunicacao proposto por NVSHMEM, no
qual a GPU descobre o destino da comunicacao e tambem emite a operacao remota.

Por outro lado, essa versao tambem introduz restricoes. Como a sincronizacao
entre blocos e feita dentro de um kernel cooperativo, a quantidade de blocos deve
ser escolhida de forma que todos possam progredir sem causar deadlock. A
implementacao precisa considerar a ocupacao do dispositivo e o numero de PEs que
compartilham a mesma GPU. Alem disso, o uso de barreiras globais no device ainda
mantem a BFS rigidamente separada por niveis; a versao somente kernel remove a
CPU do caminho entre niveis, mas nao remove a necessidade algoritmica de
sincronizar a fronteira global.

Assim, a versao somente kernel representa um passo alem da implementacao
host/device. Ela preserva a mesma semantica de descoberta por fronteiras e o
mesmo uso de atomicos remotos para definir predecessores, mas desloca a
coordenacao da BFS para dentro do kernel. Essa mudanca torna a implementacao
mais alinhada ao objetivo de usar NVSHMEM para comunicacao iniciada pela GPU e
serve como base para avaliar quanto do custo da versao anterior vinha da
participacao recorrente do host.

\bibliographystyle{plain}
\bibliography{bibliografia}
\end{document}
